%\input ../util/bookmacros.tex
%\input ../util/docparams.tex

\chapter{Integral Calculus}

\section{Fundamental Theorem of Calculus}
Keeping with the car theme, let $x(t)$ and $v(t)$ represent the distance and velocity
of a car at time $t$. 
We found a connection between these two 
functions in the last chapter: the ``instantaneous'' difference of $x(t)$ over the ``instantaneous'' difference in time is 
$x'(t) = v(t)$.
We now explore the relationship between these two functions over an {\it interval\/} of time.
Recall that a similar discrete relationship led to the \index{Fundamental Theorem of Difference Calculus}.
If $\Delta x_n = v_n$, then the sum over the ``time'' values, $0,1,\ldots,N-1$, of the difference,
$\Delta x_n$, is  $\sum_{n=0}^{N-1} \Delta x_n = \sum_{n=0}^{N-1} v_n = x_N - x_0$.
We seek an analog of this discrete result. 

The sense of the result we would expect is:
The ``continuous'' sum of $x'$ over the time interval $[t_{\rm b}, t_{\rm e}]$ is equal to
$x(t_{\rm e}) - x(t_{\rm b})$. For a start, set the discrete times 
$t_{\rm b} = t_0 < t_1 < \ldots < t_N = t_{\rm e}$ and create a {\it sequence\/}, $x$, whose $n^{\rm th}$ value, $x_n$, is defined in
terms of the {\it function\/}, $x$, by  $x_n = x(t_n)$. 
Consider the following discrete sum:\footnote{\kern 0.5pt \raise 0.5ex \hbox{\dag}}%
{The symbol $\approx$ means approximately equal.}
$$
\eqalign{
  \sum_{n=0}^{N-1} x'(t_n) \Delta t_n & \approx
  \sum_{n=0}^{N-1} {\Delta x_n \over \Delta t_n} \Delta t_n \cr
& = \sum_{n=0}^{N-1} \Delta x_n = x_{N} - x_0 = x(t_{\rm e}) - x(t_{\rm b})
}
$$ 
Note the approximation $x'(t_n) \approx {\Delta x_n \over \Delta t_n}$ just says that the derivative of $x$ at $t_n$ is
approximately the  ``rise'' between $x_n$ and $x_{n+1}$, $x_{n+1} - x_n$, divided by the ``run'', $t_{n+1} - t_n$.

We now associate a physical meaning with the above sum and extend this \index{discrete sum} to a ``continuous'' one.
We start with the simplest case: Let $v$ be the constant function having 
value $v_0$ over the time interval $[t_0, t_1]$. Then 
the change in distance 
from time $t_0$ to $t_1$ is (from high school physics) 
$x(t_1) - x(t_0) =  v_0 (t_1 - t_0)$. Another way to interpret this is that the average
velocity of the time interval $[t_0, t_1]$ is 
$v_0 = \frac{x(t_1) - x(t_0)}{t_1 - t_0}$.
Suppose instead the car velocity is a piecewise constant function with 
two values over the interval. The velocity has the constant value $v_0$ over 
the interval $[t_0, t_1)$%
\footnote{\kern 0.5pt \raise 0.5ex \hbox{\ddag}}{The half open interval $[a,b)$ is the same set as 
$[a,b]$ excluding the point  $b$.}, 
and has the constant value 
$v_1$ over the interval $[t_1, t_2]$. Forgetting for the moment how 
the car manages to jump instantly from one 
velocity to the next, 
what is the relationship 
between the distance function and the velocity function? We know how to compute 
the change in the distance function over each sub-interval, so the total 
change is the sum of the changes. 
$$
\eqalign{
x(t_2) - x(t_0) & = 
x(t_1) - x(t_0)  +  x(t_2) - x(t_1) \cr
& = \Delta x_0 + \Delta x_1 \cr
& = {\Delta x_0 \over \Delta t_0} \Delta t_0 + {\Delta x_1 \over \Delta t_1} \Delta t_1 \cr
& = v_0 \Delta t_0 + v_1 \Delta t_1 \cr
}
$$
%% KLUDGE: Had to add a break statement to avoid sill over
Suppose instead the car velocity is
a \index{piecewise constant function}, with velocities \break $v_0, v_1, \ldots, v_{N-1}$ over
intervals $[t_0, t_1), [t_1, t_2), \ldots, [t_{N-1}, t_N]$.
Again, let $x_n$ be the sequence defined by $x_n = x(t_n)$. 
What is the change in the distance over the 
interval now?
The answer is just the sum of all the changes produced by each interval 
of constant velocity. 
$$
\eqalign{
x(t_N) - x(t_0) & =
\Delta x_0 + \Delta x_1 + \ldots + \Delta x_{N-1} \cr
& = {\Delta x_0 \over \Delta t_0} \Delta t_0 + 
{\Delta x_1 \over \Delta t_1} \Delta t_1 + 
{\Delta x_2 \over \Delta t_2} \Delta t_2  
+ \cdots + {\Delta x_{N-1} \over \Delta t_{N-1}} \Delta t_{N-1} \cr
& = v_0 \Delta t_0 + v_1 \Delta t_1 + v_2 \Delta t_2 + \cdots + 
v_{N-1} \Delta t_{N-1} \cr
}
$$
This can be written in summation notation as:
$$
\eqalign{
x(t_N) - x(t_0) & = \sum_{n=0}^{N-1} v_n \Delta t_n
}
$$
Since the velocity has value $v_n$ over the interval $[t_n, t_{n+1})$, the area under that 
portion of the velocity graph is $v_n \Delta t_n$ (height $\times$ width). The area under the velocity 
graph from $t_0$ to $t_N$ is just the sum of these, which is the right hand side of the above formula.

{\it That is, the area under the velocity graph (over an interval) is the change in distance (over the same interval).\/}

In the figure below, the dotted lines are not part of the graph of the 
velocity function but have been added so that the distance traveled can be 
seen geometrically as the {\it area under the velocity graph\/}. The distance 
over the interval $[0,8]$ is the sum of the areas of each of the boxes in 
the figure below. That is, 
$x(t_8) - x(t_0) = 30 * 2 + 60 * 2 + 40 * 2 + 80 * 2 = 420$.
\epsfigure{eps/pwint.eps}{Piecewise constant velocity.}
The corresponding graph of the distance function, $x$, appears below.
\epsfigure{eps/pwintd.eps}{Piecewise linear distance function.}
We repeat what we have discovered:
No matter how complicated a piecewise constant velocity 
function is, the area under a portion of its graph represents  
the change in distance over that segment.

But what about velocity functions which aren't piecewise constant? 
How do we find the distance traveled? Well, we could approximate the velocity
function by a piecewise constant function, $v_{\rm approx}$. We will call 
the corresponding distance function $x_{\rm approx}$.
The smaller the intervals of 
constancy, the better the approximation would be to the actual velocity 
function. Specifically, we divide the interval $[t_{\rm b}, t_{\rm e}]$ 
into $N$ intervals, each with width $h_N = (t_{\rm e} - t_{\rm b}) / N$. 
Letting $t_n = t_{\rm b} + n \, h_N$, for $n = 0, 1, \ldots, N$, we define
the piecewise constant approximation function to take the value
$v(t_n)$ on the interval $[t_n, t_{n+1})$. That is, the approximation takes 
the value of the function $v$ at the left end point of a given interval and stays at this 
value for the duration of the interval. In other words, the approximation is a piecewise constant function.
For an example, see the figure below where $N$ is $10$.
\epsfigure{eps/pwcapprox.eps}{Piecewise constant approximation to velocity 
function, $v$.}
From the figure, the distance traversed by an object with the piecewise constant 
function, $v_{\rm approx}$ is as before, the sum of the velocities $\times$ the 
corresponding width of the time interval, which is $10$. That is, 
$x_{\rm approx}(100) - x_{\rm approx}(0) = v(0) * 10 + v(10) * 10 + 
v(20) * 10 + \cdots +
v(80) * 10 + v(90) * 10$.
Note that each term in the sum is the area of one of the columns in the figure. 
So, the distance traversed when the velocity profile is $v_{\rm approx}$ is the sum of all the 
boxes under the curve. 
In summation notation, we may rewrite the expression for $x_{\rm approx}$ as
$$
x_{\rm approx}(100) - x_{\rm approx}(0) = \sum_{n=0}^{9} v(10\,n) 10
$$
Here is a piecewise constant function which does a better job of 
approximating $v$.
\epsfigure{eps/pwcapprox2.eps}{A better approximation to $v$.}
The total distance traversed by an object with velocity $v_{\rm approx2}$ is again the area of its 
graph, which is the sum of all the rectangular dashed boxes. That is,
$$
x_{\rm approx2}(100) - x_{\rm approx2}(0) = \sum_{n=0}^{19} v(5\,n) 5
$$
You can see that if we make more refined 
approximations to $v$, the value of the corresponding distance function 
gets closer and closer to the area 
under the graph of $v$ and I hope you would agree that this distance function 
will get closer and closer to $x$.
Therefore, the area under the curve of $v$ is the limit of these refined
sums. This is the same {\it approximate, then take a limit\/} move we
used to define the derivative.
If we divide the interval $[0,100]$ into $N$ intervals then  the approximation would be
$\sum_{n=0}^{N-1} v(h_N \,n) h_N$ where $h_N = (100-0)/N$.
Just as when we defined derivatives, it would be useful for us
to have a notation for the limit of this process.
\goodbreak
\proclaim Definition(\index{Integral}). Given a function $f$ on an interval
$[a,b]$. For any integer $N$ we chop up the interval $[a,b]$ into $N$ uniform pieces of length $h_N = (b-a)/N$.
By the notation $\int_{a}^{b} f(t) \, dt$ we 
mean\footnote{\kern 0.5pt \raise 0.5ex \hbox{\dag}}{{\rm This definition only 
works for well behaved functions. For example, it works for 
continuous functions or a patch 
work of continuous functions. See the ``Computational Strategy for Limits'' section of the appendix 
``Infinite Sequences and Limits''  for a definition of a continuous function.}}:
$$
\int_{a}^{b} f(t) \, dt \;\; {\buildrel {\rm def} \over \equiv} \;\; 
\lim_{N \rightarrow \infty} \sum_{n=0}^{N-1} f(t_n) \Delta t_n
$$
where $t_n = a + n\, h_N$ is the left end point of the $n^{\rm th}$ interval.

The sum on the right represents an increasingly refined approximation to the area under the curve using
stepwise approximations.\footnote{\kern 0.5pt \raise 0.5ex \hbox{\ddag}}%
{The right hand side is really 
$\lim_{N \rightarrow \infty} S_N$, where $S_N \; {\buildrel {\rm def} \over \equiv} \; \sum_{n=0}^{N-1} f(t_n) \Delta t_n$.
That is, this is really the \index{limit} of a \index{sequence}, something we have already discussed.}


We say that $\int_{a}^{b} f(t) \, dt$ is the {\it integral\/} of $f$ over 
the interval $[a,b]$. We call the process of computing the 
integral {\it \index{integration}\/}.

{\bf Note:\/} The symbol, $\int$, is a stretched out `S', the first letter
of ``Sum'', while $dt$ is interpreted as a single symbol -- not multiplication of $d$ times $t$ --
and is an echo of the symbols $\Delta t_n$ used in the approximating \index{discrete sum}s.

{\bf Note:\/} We chopped $[a,b]$ into {\it uniform\/} pieces only for
simplicity; the general partitions $t_0 < t_1 < \cdots < t_N$ used earlier in
this chapter work as well, provided all the widths shrink to $0$. For the
well behaved functions of this book every such scheme produces the same
limiting value.

{\bf Note:\/} The variable $t$ (and its cousin $dt$) in the expression $\int_{a}^{b} f(t) \, dt$
are ``dummy'' variables. Their notational purpose is to represent
the ``continuous'' looping variable used in this ``sum''. Just as in programming
languages, the name is not important. It is important, however, that it
not be a name which causes confusion. For instance, if the interval
over which we are integrating is $[a,t]$ then the integral,
$\int_{a}^{t} f(t) \, dt$, doesn't make sense. This is directly analogous 
to writing a ``for loop'' that confuses the looping variable with 
a variable used to determine the end of the loop. The confused code 
might look like this in the language `C':
\break{\tt for(i = 0; i < ${\bf i\/}$; i++) $\{ \ldots \}$}.
This loop can be written correctly by using a different looping variable:
\break{\tt for(j = 0; j < i; j++) $\{ \ldots \}$}.
The incorrectly written integral, $\int_a^t f(t) \, dt$,  can be fixed in the same way:
$\int_a^t f(s) \, ds$.
To make the point clear, let us note that all of the integrals below mean 
the same thing:
$$
\int_{a}^{b} f(t) \, dt = \int_{a}^{b} f(s) \, ds = \int_{a}^{b} f(u) \, du
 = \int_{a}^{b} f(x) \, dx
$$

Our new shiny notation says that 
$\int_{t_{\rm b}}^{t_{\rm e}} v(t) \, dt = $ 
{\it the area under the graph of v}\footnote{\kern 0.5pt \raise 0.5ex \hbox{\dag}}%
{For this statement to be true we are implicitly assuming that the velocity function is non-negative.}. This is also the limit of 
$x_{\rm approx}(t_{\rm e}) - x_{\rm approx}(t_{\rm b}) $ as the approximations
get closer to $x$. Therefore, if we believe that the $x_{\rm approx}$ values 
at $t_{\rm b}$ and $t_{\rm e}$ approach the values of $x$ at 
$t_{\rm b}$ and $t_{\rm e}$, then we must have 
$$
\int_{t_{\rm b}}^{t_{\rm e}} v(t) \, dt = x(t_{\rm e}) - x(t_{\rm b})
$$
So, just as with the piecewise constant functions it is the case 
that the 
area under $v$ is just the difference in the value of the corresponding 
distance function, $x$, at the endpoints.

We can also relate the point-wise relationship of $v$ to $x$ with the 
above interval relationship. Since $x'(t) = v(t)$, we may write the above as 
$$
\int_{t_{\rm b}}^{t_{\rm e}} x'(t) \, dt = x(t_{\rm e}) - x(t_{\rm b})
$$
It turns out that this formula is true even if $x$ is not a distance function.
That is, given 
any function $F$, with corresponding derivative, $F'$, we may write 
what is called 
{\it the Fundamental Theorem of Calculus\/}:%
\proclaim Theorem(Fundamental Theorem of Calculus). If $F$ is a function such that $F'$ exists and is continuous 
over the interval $[a,b]$, then
$$
\int_{a}^{b} F'(t) \, dt = F(b) - F(a) 
$$

This is really the continuous analog of the \index{Fundamental Theorem of 
Difference Calculus}:
$$
\obrace{\sum_{n=0}^{N-1}}{Discrete Sum} 
\obrace{\Delta F_n}{Telescoping Difference} = \obrace{F_N - F_0}{Difference at End Points}
$$
We write the \index{Fundamental Theorem of Calculus} (\index{FTC}) again with suggestive markup to show the analogy with 
the above discrete equation:
$$
\obrace{\int_{a}^{b}}{Continuous Sum} 
\obrace{F'(t) \, dt}{Infinitesimal Telescoping Difference} = 
\obrace{F(b) - F(a)}{Difference at End Points}
$$
We will show how this result follows by
breaking up the interval $[a,b]$ into $N$ pieces (where $N$ is large) and approximating the integral with
a sum of \index{telescoping difference}s of a related sequence. Then we can use the Fundamental Theorem of {\it Difference\/} 
Calculus (\index{FTDC}) to show the \index{FTC}. 

To implement this program, note that each interval 
will have length $h_N = (b-a)/N$. Let $F$ be a sequence 
(we use the same name as the function) whose 
$n^{\rm th}$ element, 
$F_n$,  is the value of left 
end point of the $n^{\rm th}$ interval (intervals are labeled starting at zero). 
That is, $F_n = F(t_n)$ where $t_n$ represents the left end point of the
$n^{\rm th}$ interval; i.e., $t_n = a + n \, h_N$.
Notice that
$F_0 = F(a)$ and $F_N = F(b)$. 
In addition, note that $t_{n+1} = t_n + h_N$ ($\Delta t_n = h_N$) 
so that $F_{n+1} = F(t_{n+1}) = F(t_n + h_N)$. 
We also know that
$F'(t_n) \approx \frac{\Delta F_n}{\Delta t_n}$.%
\footnote{\kern 0.5pt \raise 0.5ex \hbox{\dag}}%
{The approximate derivative is obtained by
dividing the ``rise'' in $F$, $\Delta F_n$, by the ``run'', $\Delta t_n$. }
We now perform a crude calculation to show why the Fundamental Theorem is true:
$$
\int_{a}^{b} F'(t) \, dt \approx
\sum_{n=0}^{N-1} F'(t_n) \Delta t_n \approx 
\sum_{n=0}^{N-1} {\Delta F_n \over \Delta t_n} \Delta t_n = 
\obrace{\sum_{n=0}^{N-1} \Delta F_n = F_N - F_0}{FTDC} = F(b) - F(a)
$$
This approximation can be seen by performing a slightly more careful analysis:
$$
\eqalign{
\int_{a}^{b} F'(t) \, dt & \approx \sum_{n=0}^{N-1} F'(t_n) \Delta t_n \cr
& = \sum_{n=0}^{N-1} \bigl({\Delta F_n \over \Delta t_n}  + 
    {\rm error}(t_n)\bigr) \Delta t_n = \sum_{n=0}^{N-1} \Delta F_n  + 
\obrace{\sum_{n=0}^{N-1} {\rm error}(t_n) \Delta t_n }{very small term} \cr
& \approx  \sum_{n=0}^{N-1} \Delta F_n = F_N - F_0 = F(b) - F(a) 
}
$$
Here, ${\rm error}(t_n)$ is the correction term to make $\Delta F_n / \Delta t_n$ equal to $F'(t_n)$.
The only hard part of our crude calculation is showing that 
$\sum_{n=0}^{N-1} {\rm error}(t_n) \Delta t_n$ goes to zero as $N$ gets large.
If we call ${\rm error}_{\rm max}$ the largest of the 
${\rm error}$s 
in all the intervals, then we may bound the ``very small term'' expression above by 
${\rm error}_{\rm max} \sum_{n=0}^{N-1} \Delta t_n$.
This, in turn, is bounded by 
${\rm error}_{\rm max} (b-a)$\footnote{\kern 0.5pt \raise 0.5ex \hbox{\ddag}}%
{$\sum_{n=0}^{N-1} \Delta t_n = t_N - t_0 = b - a$ by the FTDC.}.
We need now
only believe that the maximum error term will go to zero as $N$ goes 
to infinity. 

How do we use this theorem? For instance, suppose I want to compute 
the area under the curve of a function $f$ over the interval $[a,b]$, 
how does this theorem help?
Well, we need to find an $F$ such that $F' = f$. Then 
$$
\int_a^b f(x) \, dx = \int_a^b F'(x) \, dx = F(b) - F(a)
$$
We call $F$ an {\it \index{anti-derivative}\/} of $f$. It turns out that $F$ is not 
unique; but all \index{anti-derivative}s differ by a constant.
It doesn't matter 
which one you use, you will get the same answer%
\footnote{\kern 0.5pt \raise 0.5ex \hbox{*}}{If $G$ is another anti-derivative then $G(x) = F(x)+c$ for some constant $c$.
So $G(b) - G(a) =  (F(b) + c) - (F(a) + c) = F(b) - F(a)$.}. Now we need to be 
able to find an anti-derivative of a function and we'll be done.

Let's see how to apply these ideas to a problem. Keeping with the car theme, 
suppose that you didn't have 
sensors in your car, but you were told that the velocity of the car over the 
interval $[0,10]$ (in hours) would be $v(t) =  t^2$ mph. 
How far will the car have
traveled over this time period? With our new notation we can write down the
answer $\int_0^{10} t^2 \, dt$. Although this answer is concise, the
notation hides the real work involved in its computation. That is, to compute
this value from its definition would be a great deal of work.
But, the \index{Fundamental Theorem of Calculus} says we may avoid that work 
by evaluating a certain  function 
$x$ at $10$ and $0$ and taking the difference. It is a function that has 
the property that $x'(t) = v(t) = t^2$. Can we guess what functions have 
this property? We are looking for an {\it \index{anti-derivative}\/} of $v$. 
We know that from the chapter on differentiation that  
differentiation of power functions reduces their degree by $1$. To find 
an \index{anti-derivative} we should do the opposite; that is, look for a function 
of degree one more than $2$.
Let's try $x(t) = t^3$. In this case, $x'(t) = 3t^2$. This isn't
quite right, we're off by a factor of $3$. If we let $x(t) = t^3/3$ then 
$x'(t) = t^2$ and we can find the distance traveled, it is 
$x(10) - x(0) = 10^3/3 - 0^3/3 = 1000/3 = 333 \, \frac{1}{3}$ miles.
That is,%
\footnote{\kern 0.5pt \raise 0.5ex \hbox{\dag}}%
{One last bit of notation, we use $f(t)\,\big|_{t=a}$ to mean $f(a)$. This is similar to
notation you would see in other texts.}
$$
\int_0^{10} t^2 \, dt = t^3/3\,\big|_{t=10} - t^3/3\,\big|_{t=0} = 333 \,
\frac{1}{3}
$$
We make a formal statement of the anti-derivative style FTC.

\proclaim Theorem(FTC \index{Anti-derivative} Style). If $f$ is a function and 
$F$ is an anti-derivative of $f$, then the following is true:
$$
\int_a^b f(x) \, dx = F(b) - F(a)
$$

{\bf proof:\/} This just says $\int_a^b f(x) \, dx = 
\int_a^b F'(x) \, dx = F(b) - F(a)$ 
(since $F' = f$). And the last equality is just the statement 
of the fundamental theorem of calculus.

So, in order to compute an integral of a function we need to find an 
anti-derivative of that function.

\exercise{Think about why all anti-derivatives of a function $f$ differ 
by a constant. Hint: Let $F$ and $G$ be two anti-derivatives of $f$. 
Show that the derivative of $F-G$ is zero. What kind of functions 
have zero slope at every point?}

\section{Integral Calculus: Computational Strategy}
We would like to be able to compute an {\it anti-derivative\/} of a function 
in order to be able to integrate it. We use the notation 
$\int f(x) \, dx$ to indicate an anti-derivative. That is, when we 
omit the boundaries of the integral we mean an \index{anti-derivative}; with the 
boundaries we mean the integral over an interval. 

We proceed in the same manner as 
with the \index{differential calculus}. First, we find anti-derivatives for a base 
class of functions. Next, we find formulas for anti-derivatives of functions 
composed in some way from simpler functions. Just as with the differential 
calculus, we use the power functions as a base set.

\section{A Base -- a Family of Functions and their Anti-derivatives}
\proclaim Theorem(Anti-derivative of Power Functions). An anti-derivative 
of the power 
function $f(x) = x^n$ is $F(x) = x^{n+1} / (n+1)$; valid for integer $n$, $n \neq -1$; valid for all $x$ when $n \ge 0$ and for
all $x \neq 0$ when $n < 0$. This may be written:
$$
\int x^n \, dx = x^{n+1} / (n+1)
$$
Note that when $n = -1$ we already know an anti-derivative (only valid for $x>0$): $\int 1/x \,dx = \ln(x)$.

We will need a slightly larger base than just the power functions.
There is nothing new here -- each of the formulas below is just a derivative
formula from the chapter on differentiation read in reverse. Restating the
formula for $1/x$ along with the others:
$$
\eqalign{
\int \cos(x) \, dx & = \sin(x) \cr
\int \sin(x) \, dx & = -\cos(x) \cr
\int e^x \, dx & = e^x \cr
\int 1/x \, dx & = \ln(x) \quad (x > 0) \cr
}
$$

\section{Enlarging the Class of Functions to Integrate}
Here is a result for the second stage of our program. It is a linearity 
theorem similar to the one for summation calculus.
\proclaim Theorem(Linearity of the Anti-derivative). An anti-derivative of the 
function $f$, where \allowbreak $f(x) = c_1 \, f_1(x) + c_2 \, f_2(x)$ is 
$c_1$ times an anti-derivative 
of $f_1$ plus $c_2$ times an anti-derivative of $f_2$. This may be written:
$$
\int f(x) \, dx = c_1 \int f_1(x) \, dx + c_2 \int f_2(x) \, dx
$$

Note that when $f_2(x) = 0$, the above implies $\int c f(x) \, dx = c \int f(x) \, dx$.
\proclaim Corollary(Linearity of the Anti-derivative). An anti-derivative of the 
function $f(x) = c_1 \, f_1(x) + c_2 \, f_2(x)  + \ldots + c_n f_n(x)$ is 
$$
\int f(x) \, dx = c_1 \int f_1(x) \, dx + c_2 \int f_2(x) \, dx + \ldots + c_n \int f_n(x) \, dx 
$$

This follows from the theorem on linearity and mathematical induction.
The corollary and the results for the base family can be combined to find an anti-derivative of {\it any\/} 
polynomial.

\proclaim Theorem(Anti-derivative of a Polynomial). An anti-derivative of the 
polynomial $p(x) = c_0 + c_1 x + c_2 x^2 + \cdots + c_n x^n$, ($n$ a non-negative integer) is
$$
\int p(x) \, dx = 
c_0 x + c_1 x^2/2 + c_2 x^3/3 + \cdots + c_n x^{n+1}/(n+1)
$$

\example{Compute the integral: $\int_1^2 (1 + 2x + x^2) \, dx$.}
We can use
the last theorem to find an anti-derivative of $1 + 2x + x^2$; it gives
$x + 2x^2/2 + x^3/3 = x + x^2 + x^3/3$. Therefore, by the fundamental
theorem of calculus: $\int_1^2 (1 + 2x + x^2) \, dx =
(x + x^2 + x^3/3)\,\big|_{x=2} - (x + x^2 + x^3/3)\,\big|_{x=1} = (2 + 4 + 8/3)
- (1 + 1 + 1/3) = 26/3 - 7/3 = 19/3$.


Here is another result from stage 2 of our program.
\subsection{Integration by Parts}
This is a technique which is essentially a way of computing an integral by 
using the \index{product rule} from \index{differential calculus}. The product rule states
that $(f\,g)' = f'\,g + f\,g'$. We rewrite this as: $f\, g' = (f\,g)' - f'\, g$. Integrating both sides over the interval $[a,b]$ 
gives
$$
\int_a^b f(x)g'(x) \, dx =  \int_a^b (f(x)\,g(x))' \, dx - \int_a^b f'(x)g(x) \, dx 
$$
The first integral on the right hand side is easy to compute by 
the fundamental theorem. Since this integral is easy to evaluate by the FTC, we have a direct relationship 
between the other two integrals.
$$
\int_a^b f(x)g'(x) \, dx =  \bigl(f(b)g(b) - f(a)g(a)\bigr) - \int_a^b f'(x) g(x) \, dx  
$$
We call this result {\it \index{Integration by Parts}\/}  which is the continuous analog of 
{\it \index{Summation by Parts}\/} introduced in chapter 2. 
It relates one integral to another. This is useful (just as it was for summing sequences) because if
one of the integrals can be easily evaluated we get the value of the other one for free.

\example{Compute the integral $\int_0^{\pi/2} x\sin(x) \, dx$.}
We can use the integration by parts formula here.
Notice it would be convenient 
to treat $\sin$ as the derivative of a function and then pass the derivative 
onto $x$. Let $g(x) = -\cos(x)$ and $f(x) = x$. Then $g'(x) = \sin(x)$. 
Integration by parts gives us:
$$
\eqalign{
\int_0^{\pi/2} \obrace{x\mathstrut}{$f(x)$} \obrace{\sin(x)}{$g'(x)$} \, dx &= 
\obrace{\bigl((\pi/2) * (-\cos(\pi/2)) \; - \; 0  * (-\cos(0))\bigr)}{$f(\pi/2)  g(\pi/2) \;-\; f(0)  g(0)$} - \int_0^{\pi/2} \obrace{x'}{$f'(x)$} \obrace{-\cos(x)}{$g(x)$} \, dx   \cr
&= - \int_0^{\pi/2} -\cos(x)\, dx =  \int_0^{\pi/2} \cos(x)\, dx = \sin(\pi/2) \;-\; \sin(0) \cr
&= 1
}
$$
Note that the values $0$ and $\pi/2$ are in radians. The value $\pi/2$ radians is 90 degrees. This means that
$\sin(\pi/2) = 1$ and $\cos(\pi/2) = 0$.
\subsection{Substitution}
Consider the integral $\int \sin^2(x) \cos(x) \, dx$. None of the rules we have
developed so far will crack it -- the integrand is not a polynomial, nor is it a
combination of our base functions that linearity or integration by parts can untangle.
We need one more technique. It is the {\it change the representation\/} move
of the preface: rewrite the integrand until a form we already know how to
integrate appears.
This technique uses the \index{chain rule} (from \index{differential calculus}) to compute
an integral.
The chain rule states that $f(u(x))' = f'(u(x)) u'(x)$. Integrating both 
sides gives
$$
\int_a^b f'(u(x)) u'(x) \, dx = \int_a^b f(u(x))' \, dx = f(u(b)) - f(u(a))
$$
Successful use of the above techniques is based on how well one can pattern match.
It takes practice, but it is not our focus here. We demonstrate the idea through a few examples.

\example{Compute $\int_a^b (x-a)^n\, dx$ when $n$ is a non-negative integer.}
Letting $u(x) = (x-a)$, we see that the integral can be written as
$\int_a^b u(x)^n u'(x) \, dx$.
In this case letting $f'(z) = z^n$ (which has $z^{n+1}/(n+1)$ as an anti-derivative), the 
substitution formula gives
$\int_a^b f'(u(x)) u'(x)\, dx = f(u(b)) - f(u(a)) = {1 \over n+1} (b-a)^{n+1} - {1 \over n+1} (a-a)^{n+1} = {1 \over n+1} (b-a)^{n+1}$.
\example{Compute $\int_a^b (a-x)^n\, dx$ when $n$ is a non-negative integer.}
Letting $u(x) = (a-x)$, we see that the integral can be written as
$-\int_a^b u(x)^n u'(x) \, dx$.
Setting $f'(z) = z^n$ as in the previous example the 
substitution formula gives
$-\int_a^b f'(u(x)) u'(x)\, dx = -( f(u(b)) - f(u(a)) ) = -\bigl({1 \over n+1} (a-b)^{n+1} - {1 \over n+1} (a-a)^{n+1}\bigr) = -{1 \over n+1} (a-b)^{n+1}$.
\example{Compute $\int_0^1 \sin^2(x) \cos(x) \, dx$.}
Letting $u(x) = \sin(x)$, 
and therefore $u'(x) = \cos(x)$,
we see that the integral can be written as $\int_0^1 u(x)^2 u'(x) \, dx$.
In this case $f'(z) = z^2$. Using the substitution formula gives 
$\int_0^1 f'(u(x)) u'(x) \, dx = f(u(1)) - f(u(0)) = u(1)^3/3 - u(0)^3/3 =
\sin(x)^3/3\,\big|_{x=1} - \sin(x)^3/3\,\big|_{x=0} = \sin^3(1)/3 - 0 = \sin^3(1)/3$. Note, as always
the angle measure is in radians. So, the $1$ in our answer, $\sin^3(1)/3$, means $1$ radian,
which is approximately 57 degrees.

\exercise{Compute the following integrals: (1) $\int_0^1 (1 + x)^7 \, dx$;
(2) $\int_0^{\pi/2} x \cos(x) \, dx$; (3) $\int_0^1 x \, e^{x^2} \, dx$.
Hint: the first and last yield to substitution; the middle one to
integration by parts.}

\exercise{Compute the area under the graph of $\sin$ over the interval
$[0, \pi]$. Then compute $\int_0^{2\pi} \sin(x) \, dx$. Why is the second
answer $0$ even though the graph over $[0, 2\pi]$ certainly encloses area?
(Recall the footnote earlier in this chapter: the area interpretation
assumed the function was non-negative.)}

\section{Can't always find Formulas for Anti-derivatives}
As noted in the last chapter, not every function has a derivative. And even when it does, it may not 
be easy to find a formula for the derivative. However, if the function has a nice formula it is often the 
case that we can find
a formula for the derivative. This is not the case with anti-derivatives. For nice continuous functions we always 
have an \index{anti-derivative}, but finding a formula is not always possible -- even for relatively simple functions. 
That is, there are functions with 
nice formulas for which there is no formula for an anti-derivative. One example is the bell shaped 
functions associated with the subjects of probability and statistics. 
The graph of the function $f(x) = e^{-x^2}$ is bell shaped and its formula is simple.
The area under the graph of this function exists and  we can {\it create\/} a function which is the area 
under this curve from $0$ to $x$.
We can write the value down in our notation as: $A(x) = \int_0^x e^{-s^2} \, ds$. 
However, there is no formula for this function in terms of powers of 
$x$, logarithmic, exponential, or trig functions because there is no formula for an anti-derivative of $e^{-x^2}$. 
The integral (the area under the graph)
exists; that is, the function $A(x)$ is a valid function -- there is just no formula for it. See the 
appendix ``Functions, Formulas, Graphs, and Notation'' for other examples of functions without formulas.
 
\section{Summary}
The \index{fundamental theorem of calculus} can be stated succinctly:  
the integration operation is the 
inverse of the \index{differentiation operation}. This can be seen from the \index{FTC}:
$\int_{a}^b f'(t) \, dt = f(b) - f(a)$, and writing it more crudely as $\int f' = f$. What makes 
this result useful is if for a given integral, $\int_{a}^b v(t)\, dt$, one can recognize $v$ as
the derivative of another function $f$, then this integral is 
$\int_{a}^b f'(t)\, dt$ and by the FTC this is just $f(b) - f(a)$.

The process of recognition is helped greatly if there is a well developed stock of differentiation formulas. 
Then when confronted with integrating a function, we may be able to recognize, or more easily deduce, 
a function which
is an {\it anti-derivative\/} and then easily compute the integral. As with differencing and summing, finding
formulas for derivatives and \index{anti-derivative}s of functions is much like multiplying and dividing numbers. 
Computing derivatives is relatively easier -- like multiplying, while finding 
anti-derivatives is harder -- like division. To 
do long division it helps to be fluent with multiplication, because trial and error multiplications are
necessary. The same is true when computing anti-derivatives, it helps first to be fluent with the formulas for
the derivatives of a number of functions. To this end, one would
\beginEnum
\enum{Develop formulas for the derivatives of a base collection of functions.}
\enum{Develop formulas for the derivatives of functions formed in various ways from the base.}
\endEnum

This is exactly what we did in the last chapter.

Stepping back, this whole chapter was really the discrete game replayed on the
continuous board -- every move we made had already been rehearsed with sequences,
differences, and sums. Here is the dictionary:
\bigskip
\centerline{\vbox{\halign{\quad\hfil#\hfil\quad&\quad\hfil#\hfil\quad\cr
{\bf Discrete}&{\bf Continuous}\cr
\noalign{\smallskip\hrule\smallskip}
$\Delta$&$d/dx$\cr
$\sum$&$\int$\cr
Fundamental Theorem of Difference Calculus (telescoping)&Fundamental Theorem of Calculus\cr
Summation by Parts&Integration by Parts\cr
anti-difference&anti-derivative\cr
}}}
\bigskip


